% !TeX program = pdflatex
% DLR-style cover letter for JFM resubmission (JFM-2026-1781 referenced).
% Body rewritten 2026-09-09 (task U4/M6): objection-by-objection discharge
% + one-paper rationale.  Template/header unchanged from the July letter.
%
% Put dlrlogo.png in the same folder as this .tex file.

\documentclass[11pt,a4paper]{article}

\usepackage[a4paper,left=25mm,right=25mm,top=22mm,bottom=22mm]{geometry}
\usepackage{graphicx}
\usepackage[absolute,overlay]{textpos}
\usepackage{xcolor}
\usepackage{helvet}
\usepackage[T1]{fontenc}
\usepackage[utf8]{inputenc}
\usepackage[english]{babel}
\usepackage{setspace}
\usepackage{parskip}
\usepackage{ragged2e}
\usepackage{microtype}
\usepackage[hidelinks]{hyperref}

\renewcommand{\familydefault}{\sfdefault}
\pagestyle{empty}
\setlength{\parindent}{0pt}
\setstretch{1.02}

% -----------------------------------------------------------------------------
% Editable letter data
% -----------------------------------------------------------------------------
\newcommand{\Institute}{Institut für Aerodynamik und Strömungstechnik}
\newcommand{\DLRAddress}{Lilienthalplatz 7, 38108 Braunschweig}

\newcommand{\Recipient}{%
The Editor\\
Journal of Fluid Mechanics\\
Cambridge University Press%
}

\newcommand{\ContactPerson}{Dr. Sparsh Sharma}
\newcommand{\Phone}{0531 295--1530}
\newcommand{\Email}{sparsh.sharma@dlr.de}
\newcommand{\LetterDate}{\today}

\newcommand{\Subject}{Submission of a reconstructed manuscript
(ref.\ JFM-2026-1781) to the Journal of Fluid Mechanics}

\newcommand{\ManuscriptTitle}{Strain-coupled one-dimensional turbulence:
a single-line closure for rapid distortion, its measured limits, and the
consequence for leading-edge noise}

\newcommand{\ClosingName}{Dr. Sparsh Sharma}

% Logo file
\newcommand{\LogoFile}{dlrlogo.png}

% -----------------------------------------------------------------------------
% Small helpers
% -----------------------------------------------------------------------------
\newcommand{\ContactLabel}[1]{\textcolor{black!55}{\scriptsize #1}}

\begin{document}

% -----------------------------------------------------------------------------
% Header
% -----------------------------------------------------------------------------
\begin{textblock*}{95mm}(25mm,28mm)
  {\small \Institute}
\end{textblock*}

\begin{textblock*}{75mm}(126mm,20mm)
  \raggedleft
  \includegraphics[width=58mm]{\LogoFile}
\end{textblock*}

% -----------------------------------------------------------------------------
% Recipient address block and contact block
% -----------------------------------------------------------------------------
\begin{textblock*}{86mm}(25mm,66mm)
  {\scriptsize
  DLR e.\,V.\hspace{3mm} \Institute\\[-0.2mm]
  \hspace*{21mm}\DLRAddress}
  \vspace{7mm}

  {\normalsize \Recipient}
\end{textblock*}

\begin{textblock*}{65mm}(126mm,66mm)
  \raggedleft
  \begin{tabular}{@{}r@{\hspace{7mm}}l@{}}
    \ContactLabel{Reference}       & JFM-2026-1781 \\
    \ContactLabel{Manuscript type} & Regular article \\
    \\[-1mm]
    \ContactLabel{Contact person}  & \ContactPerson \\
    \\[-1mm]
    \ContactLabel{Telephone}       & \Phone \\
    \\[-1mm]
    \ContactLabel{E-mail}          & \Email \\
    \\[2mm]
                                    & \LetterDate
  \end{tabular}
\end{textblock*}

% -----------------------------------------------------------------------------
% Body
% -----------------------------------------------------------------------------
\vspace*{118mm}

{\bfseries \Subject}

\vspace{11mm}

Dear Editor,

\vspace{5mm}

We are pleased to submit the manuscript entitled
``\ManuscriptTitle'' by Sparsh Sharma and Alan R.\ Kerstein for
consideration as a regular article in the \textit{Journal of Fluid
Mechanics}.

\vspace{5mm}

This manuscript draws on material from JFM-2026-1781, which was
declined on 18 August 2026, and per the instruction in that decision we
cite the number here and request that the submission be handled by the
same associate editor, Professor Sutanu Sarkar. We are not resubmitting
that paper. The present manuscript is a reconstruction around what the
referees identified as its strongest element---the single-line closure
analysis, which two referees independently singled out and one
explicitly suggested rebuilding the study around---and it contains a
body of work that did not exist at the original submission: an
independent DNS benchmark, a corrected linear reference, a measured
closure assumption, a new relaxation mechanism developed with the
model's originator (now a co-author), and the application carried
through to radiated levels. We summarise below how each of the
referees' objections has been discharged, because we took every one of
them as a work order rather than a disagreement.

\vspace{5mm}

\textbf{Validation against the model itself (Refs.\ 1.1, 2, 3).} The
central spectral results are now benchmarked against a strained-box
direct numerical simulation with exact rapid-distortion companion
calculations, all reduced to the same one-dimensional observables, so
that model, DNS and exact linear theory meet in a three-way (and, for
the relaxation mechanism, five-way) comparison on a single benchmark
observable. The strain-on/strain-off construction survives only as an
internal diagnostic; every headline claim now has an external referent.

\vspace{5mm}

\textbf{The rigid-translation ``linear reference'' (Ref.\ 3.1).} The
referee was right, and the point is now a result of the paper rather
than an error in it: we compute exact rapid-distortion theory projected
onto the line and show that the projected spectra are redistributed
broadband---by a factor of four across the resolved range at unit total
strain---so the correct linear reference differs from rigid translation
already at the linear level. Every comparison in the paper is built on
the corrected reference, and the earlier claim is withdrawn.

\vspace{5mm}

\textbf{The single-line closure ``not closed'', and not exercised
(Refs.\ 3.2, 1.2).} The closure analysis is rebuilt in exactly the
modest form the objection demands: the obstruction to a single-line
closure is stated exactly, sharp bounds on the rapid term are derived
from line data alone (the previously underived relations now derived in
an appendix), and the axisymmetry hypothesis---undefined in the
original submission---is defined and \emph{measured} rather than
assumed: exact for axisymmetric distortion, with its plane-strain
violation quantified as a function of accumulated strain. This
machinery is used throughout the paper (the validity envelope, the
deficiency localisation, the uncertainty carried into the application);
the full wavenumber-dependent kernel it defines is identified, on
measured grounds, as the remaining step, and we say so plainly rather
than claiming it.

\vspace{5mm}

\textbf{The plane-strain behaviour of the closure (Ref.\ 1.3).} The
moment-level validation against Lee \& Reynolds has been recomputed
from the exact solutions (correcting, in the process, an error in our
own earlier exact curve), and the one qualitative failure---the
neutral-direction anisotropy---is traced to the orientation information
that every single-point closure discards, with its consequence for the
spectral results assessed through the measured axisymmetry violation
rather than left unexamined.

\vspace{5mm}

\textbf{The application not carried through (Refs.\ 1.4, 2).} The
manuscript now contains the mapping the referees asked for: the model's
line spectra enter Amiet's theory through a fitted two-parameter
distorted-spectrum representation whose residual is the quoted
uncertainty, yielding both a geometry-free change-of-kernel result and
an assembled absolute sound-pressure-level prediction at the reference
operating configuration. The computed distortion lowers the predicted
level by 5--20\,dB relative to the frozen von K\'arm\'an input of
standard practice---a term too large to neglect, and of opposite sign
to the surface-blocking correction, which we state as the scope
boundary.

\vspace{5mm}

\textbf{Presentation (Refs.\ 2, 3.3).} The manuscript was rewritten
from a fresh skeleton, not edited: established material is separated
from new (the model review is explicitly marked as containing nothing
original and compressed to citations), Sagaut \& Cambon is adopted as
the rapid-distortion spine, the leading-edge distortion literature the
referees named (dos Santos, Ribeiro, Piccolo) is engaged, and the main
text is shorter than the rejected version while carrying substantially
more content; the validation against Chen, Meneveau \& Katz and the
complete intervention family are provided as supplementary material.

\vspace{5mm}

One structural decision deserves a sentence. The new results---the
closure analysis, the mechanism study and the noise consequence---could
have been split into separate papers; we judged that separated they
would each lose the evidential chain that gives them force. They are
therefore presented as one paper in three parts, with the supplementary
material absorbing what the arc does not need.

\vspace{5mm}

The manuscript is original, has not been published previously, and is
not under consideration elsewhere. Both authors have approved the
submission and take responsibility for the content of the manuscript. A
declaration of the use of AI tools is included in the manuscript in
accordance with Cambridge University Press guidance. To the best of our
knowledge, there are no conflicts of interest that would affect the
review of this work. We are of course content for the original referees
to see how their objections were used.

\vspace{5mm}

Thank you for considering this manuscript for publication in the
\textit{Journal of Fluid Mechanics}.

\vspace{10mm}

Yours sincerely,

\vspace{18mm}

\ClosingName, on behalf of the authors\\
German Aerospace Center (DLR)\\
\Institute\\
Lilienthalplatz 7\\
38108 Braunschweig, Germany\\
\href{mailto:\Email}{\Email}

% -----------------------------------------------------------------------------
% Footer
% -----------------------------------------------------------------------------
\begin{textblock*}{65mm}(25mm,260mm)
  \scriptsize
  Deutsches Zentrum für Luft- und Raumfahrt e.\,V.\\
  Institut für Aerodynamik und Strömungstechnik\\
  Lilienthalplatz 7\\
  38108 Braunschweig
\end{textblock*}

\end{document}
